\section{Legendre Polynomials}
The Legendre polynomials are the eigenfunctions of the singular \eqref{eq:Sturm-Liouville} with $p(x) = (1-x^2)$, $q = 0$, and $w = 1$
\begin{equation*}
    -\frac{d}{dx} \left( (1-x^2) \frac{dL_k}{dx} \right) = k(k+1) L_k(x).
\end{equation*}
We have a unique solution if we set the condition $L_k(1) = 1$ and it is given by
\begin{equation*}
    L_{k}(x) = \frac 1 {2^k} \sum_{l=0}^{\lfloor k / 2 \rfloor} (-1)^l \binom{k}{l} \binom{2k-2l}{k} x^{k-2l} .
\end{equation*}

\begin{example}[Legendre-Gauss integration]
    Due to \Cref{rem:Sturm-Liouville degenerate} we have that $L_k$ are orthogonal in $L^2(-1,1)$. Therefore, we can apply \Cref{thm:Gauss integration} with the precise value
    \begin{equation*}
        x_j \text{ zeros of } L_{N+1}, \qquad w_j = \frac{2}{(1-x_j^2)(\frac{dL_{N+1}}{dx}(x_j))^2}.
    \end{equation*}
\end{example}

\begin{example}[Legendre-Gauss-Lobatto integration]
    Instead of taking the extrema of $L_N$ we take
    \begin{equation*}
        x_0 = -1, x_N = 1, x_1, \cdots, x_{N-1} \text{ the zeros of } L_N'
    \end{equation*}
    The weights
    \begin{equation*}
        w_j = \frac{2}{N(N+1)} \frac{1}{(L_N(x_j))^2}
    \end{equation*}
    Based on these nodes we construct the characteristic Lagrange polynomials
    \begin{equation*}
        \psi_j (x) = \frac{1}{N(N+1)} \frac{1-x^2}{x_j - x} \frac{L_N'(x)}{L_N(x)}
    \end{equation*}
    which are discrete (shifted) delta functions, since they satisfy $\psi_j (x_j) = \delta_{jk}$. We will use them.
\end{example}

Again, we have the estimate for $k \le m$
\begin{equation*}
    \|u - P_N u\|_{H^k(-1,1)} \le C N^{k-m} \|u \|_{H^m (-1,1)}.
\end{equation*}
We also have that
\begin{equation*}
    \|u - I_N u \|_{L^2(-1,1)} \le C N^{-m} \|u \|_{H^m (-1,1)}.
\end{equation*}

\begin{example}[Galerkin method with Legendre polynomials]
    Let us try to solve the problem
    \begin{equation*}
        -\frac{d}{dx}\left((1-x^2) \frac{du}{dx}\right) = f \text{ in } (-1,1).
    \end{equation*}
    Due to \Cref{rem:Sturm-Liouville degenerate} we do not impose any boundary conditions.
    Notice that $L_k$ are orthogonal in $L^2(-1,1)$ and $\deg L_k = k$. Therefore, they generate $L^2(-1,1)$. Let us look for a solution
    \begin{equation*}
        u = \sum_{k=0}^\infty \hat u_n L_k
    \end{equation*}
    Multiplying the equation by $L_k$ and integrating by parts we deduce
    \begin{equation*}
        k(k+1) (u, L_k)_{L^2} = (f, L_k).
    \end{equation*}
    Therefore, we can write the solution as
    \begin{equation*}
        u = \sum_{k=0}^\infty \frac{(f, L_k)_{L^2}}{k(k+1) \|L_k\|_{L^2}^2} L_k.
    \end{equation*}

    Galerkin's method would suggest to look for an approximate solution in $V_N = \{ \sum_{k=0}^N a_k L_k : a_k \in \mathbb R\}$ by finding
    \begin{equation*}
        u_N \in V_N \text{ such that } \left( -\frac{d}{dx}\left((1-x^2) \frac{du_N}{dx}\right), v_N \right)_{L^2} = (f, v_N)_{L^2}, \qquad \forall v_N \in V_N.
    \end{equation*}
    But notice that then
    \begin{equation*}
        u_N = \sum_{k=0}^N \frac{(f, L_k)_{L^2}}{k(k+1) \|L_k\|_{L^2}^2} L_k = P_N u.
    \end{equation*}
\end{example}